% Ported from sections/04_introduction.md (v11, APPROVED 2026-09-18).
% Continuous prose, no subsections. Final-manuscript Sec.~I.
\section{Introduction}
\label{sec:intro}

MgO-based magnetic tunnel junctions reach very large tunnel magnetoresistance at room temperature
when the barrier is crystalline and the metal electrodes are matched to it in-plane
\cite{yuasa2004,parkin2004}. The transport calculations that account for the effect were carried out
for an ideal Fe$|$MgO$|$Fe(001) interface, in which a flat metal layer sits directly above the
substrate oxygen \cite{butler2001}. Achieving those values in practice is a question of realising
that interface closely enough: the high magnetoresistance is obtained when the barrier is
(001)-oriented and the metal film flat and continuous \cite{yuasa2004}, and the barrier's
orientation and stress relaxation are themselves sensitive to how the metal is deposited and
annealed \cite{djayaprawira2005,ikeda2008}. The structural quality of the metal/oxide interface is
therefore part of the specification of the device rather than an incidental detail.

Two recent trends in junction fabrication have only sharpened that point. The first is the push on
the tunnel magnetoresistance itself: after a decade of stagnation at the 604\,\% reported in 2008
for junctions with boron-alloy electrodes \cite{ikeda2008}, epitaxial junctions now reach a record
631\,\% at room temperature (1143\,\% at 10~K), and the gain comes explicitly from the
interface---from tuning the crystallographic orientation and the oxygen content of the MgO barrier
through ultrathin metal insertions at the metal/oxide boundary \cite{scheike2023}. The second is the
move to perpendicularly magnetised junctions for dense, non-volatile memory, in which the
metal/oxide interface is itself a functional layer: the Fe/MgO interface carries a large
perpendicular surface anisotropy, whose strength is traced to the hybridisation of interfacial iron
and oxygen states \cite{solano2022}, and sub-nanometre metal films grown on the barrier are
engineered for the perpendicular anisotropy and the low magnetic damping that low-power switching
requires \cite{ichinose2025}. Both trends converge on the same requirement---a metal film that is
flat, epitaxial and well-wetting on the crystalline barrier---and, in practice, on layer-by-layer
control of the roughness and thickness of each film in the stack \cite{ghemes2024}.

The flat, lattice-matched film assumed in such interface models is not, however, the configuration
that a metal film necessarily adopts when it is deposited on MgO(001). At the monolayer coverages
relevant to a thin electrode, growth experiments find the film dewetting into three-dimensional
islands: He-atom scattering records three-dimensional island growth of Fe on MgO(001) at room
temperature, suppressed only when the film is deposited at 140~K \cite{fahsold2000}; scanning
tunnelling microscopy combined with X-ray magnetic circular dichroism finds sub-nanometre Fe growing
three-dimensionally, with a two-dimensional growth mode reported only above about 6.5 monolayers
\cite{torelli2009}; and grazing-incidence small-angle X-ray scattering on five monolayers likewise
identifies Volmer--Weber growth with spherical islands \cite{reitinger2007}. A flat, pseudomorphic
monolayer is obtained instead by slow deposition onto a cleaved, oxygen-annealed crystal
\cite{urano1988}. The islanding tendency survives even into the most controlled modern fabrication:
grain-to-grain epitaxial growth of a sub-nanometre ferromagnetic layer on a polycrystalline
MgO(001) barrier on 300\,mm wafers must be performed at cryogenic temperature (100~K), because
room-temperature deposition lets the metal nucleate as islands and break the film's continuity
\cite{ichinose2025}. Two structural outcomes thus compete at this coverage---a flat, well-wetting
film and a dewetted island---and the growth experiments constrain their kinetics rather than their
relative energies, while interface models assume the flat geometry by construction. What is missing
is a comparison of the two configurations' energies for one and the same interface, and an answer to
how that comparison responds to a change in the film's composition.

The composition is not arbitrary. The metal electrodes of the highest-magnetoresistance MgO
junctions are boron-bearing and amorphous as deposited, which is also the condition under which the
MgO barrier grows (001)-textured \cite{djayaprawira2005}. Boron is a glass-forming addition, and
added species are in general understood to stabilise disordered configurations at the expense of
crystalline ones, in the core of the confusion principle of metallic-glass formation
\cite{greer1993}. Applied to this interface, where the flat film is the two-dimensional,
disordered-like configuration and the island the three-dimensional, ordered-like one, it suggests
that boron should lower the energy of the flat configuration relative to the island. Whether it does
so, and by how much, is a question for calculation.

Testing this idea directly is not a single-calculation question. The flat configuration is one
well-defined structure, but the ordered island it competes with is not: the dewetted film comprises
a large family of distinct structures with many local minima, so no one geometry can represent it,
and the two phases cannot be compared by relaxing a single candidate of each.
The confusion-principle question---whether an added element lowers the disordered configuration
relative to the ordered one---is therefore about the landscape as a whole, and one that is expensive
to answer, because exploring it with density-functional theory alone would require evaluating far too
many structures. Two features of the problem make it tractable. First, the question is specifically
about two phases---the flat, well-wetting film and the dewetted island---so rather than surveying
every configuration we restrict the exploration to these two targets, a deliberately biased
strategy. Second, because the search must still sample many candidates to populate both phases, we
guide it with a surrogate model, so that the expensive density-functional evaluations are spent only
where they are most informative.

In this work, by implementing a surrogate-driven active-learning search with a deliberately biased
exploration strategy \cite{gofee2017,hamamoto2023}, we categorize the potential-energy surface of Fe
on MgO(001) into its two growth modes---the flat, well-wetting film and the dewetted island---for
the same host with and without added boron, and quantify how the added element shifts the relative
energy of the flat mode.

This paper is organized as follows. In Sec.~\ref{sec:methods}, we present the active-learning
search, the two interface models, and the metrics used to separate the two growth modes. In
Sec.~\ref{sec:results}, we first establish the two-phase landscape and the island ground state, and
then report the effect of boron on the flat--island separation. In Sec.~\ref{sec:discussion}, we
discuss the origin of the island ground state and the role of the added element. We summarize our
work in Sec.~\ref{sec:conclusion}.
